
% \section{Generated Examples}

% \label{sec:app-example}

% Table~\ref{tab:example} shows errored examples.

% \section{Result with Different Retrieving Method}
% \label{sec:appendix-result-with-DiffRetrMethod}
% \input{table_ret}
% To assess the effectiveness of our pronunciation-based retrieving method, we compared it with the semantic similarity-based retriever. For text embedding, we utilized the pre-trained BERT \cite{bert} (Table \ref{tab:ret}). In our experimental results, we observed contrasting outcomes depending on the model. The Llama2 13 B model showed overwhelmingly superior performance when utilizing a retriever based on pronunciation. However, for the OPT model, examples with semantic similarity yielded better results. This implies that, although the pronunciation approach demonstrated better robustness compared to TTS-ASR or FT methods, there is still room for improvement by enhancing the retrieving method.


\section{Details of the EPA Method}






\subsection{Prompt Used for EPA}
\label{app:prompt}
The prompts used in Step 1 and Step 2 are provided below.


\begin{prompt}[colback=black!0!white, colframe=black!98!black]
{Step 1 Prompt}{\textcolor{myBlue}{Generate ASR error augmented text with similar pronunciation but different words based on the given gold text examples.}\\
Apply character and word substitutions, additions, or deletions while maintaining the overall pronunciation and context. \\
Error rate should be high 
\par            % ← end the paragraph
\noindent\dotfill\par 
\textbf{Example 1}\\
Original: they have a single naupliar eye  \\
ASR-errored: they have a single nor pure eye \\

\textbf{Example 2}\\
Original: i must have saint louis then huzza  \\
ASR-errored: i must have st louis then hazard \\

\textbf{Example 3}\\
Original: i wonder uncle did not have her come  \\
ASR-errored: i wonder uncle did not have a problem 
\par            % ← end the paragraph
\noindent\dotfill\par 
\textcolor{myBlue}{Now, following the above examples, generate an ASR-errored version of the following sentence:\\}
Original: \textcolor{myOrange}{[Target utterance]} \\
ASR-errored: 
}
\end{prompt}


\begin{prompt}[colback=black!0!white, colframe=black!98!black]
{Step 2 Prompt}{\textcolor{myBlue}{Change the key words in <hl> tag, to having a ASR error.
ASR error has similar pronounciation with the correct word, but different charater. }\\
Here is some example.
\par            % ← end the paragraph
\noindent\dotfill\par 
\textbf{Example 1}\\
Original: I want to buy a book about <hl>luwombo best</hl> restaurant. \\
Keywords : luwombo best \\
Result: I want to buy a book about luwambo vest restaurant.  \\\\
\textbf{Example 2}\\
Original: hi, i'm looking for a bus that is depart from <hl>eliot<hl/> and arriving to <hl>holiday inn williamsport<hl/>? \\
Keywords : eliot, holiday inn williamsport \\
Result:  hi, i'm looking for a bus that is depart from Ellyot and arriving to holliday inn william's port \\\\
\textbf{Example 3}\\
Original: the <hl>chabuton ramen<hl/> is a restaurant on the east. \\
Keywords : chabuton ramen \\
Result: the shabuton raymond is a restaurant on the east.\\ 
\par            % ← end the paragraph
\noindent\dotfill\par 
\textcolor{myBlue}{Now, following the above examples, generate an ASR-errored version of the following sentence:\\}
Original: \textcolor{myOrange}{[Target utterance with <hl> tag]} \\
ASR-errored: 
}
\end{prompt}


% \normalsize
% \textbf{Step 2 Prompt Example} \\
% \small{\texttt{Change the key words in <hl> tag, to having a ASR error.
% ASR error has similar pronounciation with the correct word, but different charater. \\
% Here is some example \\
% \{Example 1\}\\
% Original: I want to buy a book about <hl>luwombo best</hl> restaurant. Keywords : luwombo best \\
% Result: I want to buy a book about luwambo vest restaurant.  \\
% \{Example 2\}\\
% Original: hi, i'm looking for a bus that is depart from <hl>eliot<hl/> and arriving to <hl>holiday inn williamsport<hl/>? Keywords : eliot, holiday inn williamsport \\
% Result:  hi, i'm looking for a bus that is depart from Ellyot and arriving to holliday inn william's port \\
% \{Example 3\}\\
% Original: the <hl>chabuton ramen<hl/> is a restaurant on the east. Keywords : chabuton ramen \\
% Result: the shabuton raymond is a restaurant on the east.\\ 
% \{Example 4\}\\
% Original: i was hoping to see local faces while in cambridge. some <hl>entertainment</hl> would grate. Keywords: entertainment \\
% Result: 
% }}
% \\
% \noindent




\subsection{Detailed Keyword Highlighting Strategy for EPA}



% Below is an example for generatig target utterance with <hl> tag. \\ 
% Original: \\
% \hspace*{1em} Hi, I need to go to Green Day hotel, then book a table at the Grill House. \\
% Dialogue State: \\
% \hspace*{1em} \texttt{\{"hotel-name" : "Green Day", "restaurant-name" : "Grill House"\}}
% With \texttt{<hl>} tags: \\
% \hspace*{1em} Hi, I need to go to \texttt{<hl>}Green Day\texttt{</hl>} hotel, then book a table at the \texttt{<hl>}Grill House\texttt{</hl>}.




\begin{table}[h!]
\small
\resizebox{\columnwidth}{!}{
\begin{tabular}{p{0.2\columnwidth} p{0.8\columnwidth}}
\toprule
\multicolumn{2}{l}{\textbf{Example of adding <hl> tag}} \\
\midrule
Original & Hi, I need to go to Green Day hotel, then book a table at the Grill House. \\
Dialogue State & hotel-name: \textcolor{myGreen}{Green Day}, restaurant-name:\textcolor{myBlue}{Grill House} \\
\cdashline{1-2}
With \texttt{<hl>} tags & Hi, I need to go to \texttt{<hl>}\textcolor{myGreen}{Green Day}\texttt{</hl>} hotel, then book a table at the \texttt{<hl>}\textcolor{myBlue}{Grill House}\texttt{</hl>}. \\
\bottomrule
\end{tabular}
}
\caption{Example of keyword highlighting using \texttt{<hl>} tags based on dialogue state annotations.}
\label{tab:hl_example}
\end{table}

To explicitly introduce keyword-specific ASR errors, we first identify dialogue state values from the training corpus and match them against the user utterance ($\dot{u}$). Matched values are then automatically wrapped with \texttt{<hl>} tags based on slot annotations (e.g., DST slot labels or NER tags), as shown in Table~\ref{tab:hl_example}. These highlighted utterances are passed to the LLM, which is instructed to perturb the text within the \texttt{<hl>} tags while preserving the rest. This keyword highlighting strategy is task-agnostic and can be easily applied to other keyword-sensitive tasks such as Named Entity Recognition (NER) or Spoken Language Understanding (SLU), where certain slot values or entities are critical for downstream prediction.


\begin{table*}[h!]
\small
\resizebox{\textwidth}{!}{
\begin{tabular}{p{0.2\textwidth}p{0.8\textwidth}}


\toprule
\multicolumn{2}{l}{Example 1}     \\ \toprule
                                            
Target & \textcolor{myOrange}{can you tell} me the address to the \textcolor{myRed}{police station} \textcolor{myGreen}{in point pleasant?}  \\ \cdashline{1-2}

Retrieved 1 &  \textcolor{myRed}{frayser station} was not the depot \textcolor{myGreen}{on the point}      \\
\quad \small{+ ASR} & freya station was not the deep watch on the point   \\

Retrieved 2&  \textcolor{myOrange}{can you get me} the maldeamores saga    \\
\quad \small{+ ASR} & can you get me the melamorphos     \\
Retrieved 3               &  \textcolor{myOrange}{cannot you tell} her whom i am eh joseph\\
\quad \small{+ ASR} & cannot you tell her whom i am    \\

\toprule
\multicolumn{2}{l}{Example 2}     \\ \toprule
                                            
Target & no, i just \textcolor{myBlue}{need} to make \textcolor{myGreen}{sure} it's cheap. oh, and i need parking. \\ \cdashline{1-2}

Retrieved 1 &  i \textcolor{myBlue}{need} fifty ten foot long \textcolor{myGreen}{segments} of wire    \\
\quad \small{+ ASR} & i need fifty ten foot long signals of my life   \\

Retrieved 2&  a drive with a different encoding mechanism would \textcolor{myBlue}{need} different patterns  \\
\quad \small{+ ASR} & and drive was a different building recognition would need different patterns    \\
Retrieved 3               &  to  \textcolor{myBlue}{reach} to calcutta you  \textcolor{myBlue}{need} less time to reach dhaka   \\
\quad \small{+ ASR} & to reach tukaukara you need last time to reach daka    \\




\toprule
\multicolumn{2}{l}{Example 3}     \\ \toprule
                                            
Target & i'm \textcolor{myBlue}{open to} any \textcolor{myOrange}{kind of} food. i'm looking for something in the \textcolor{myGreen}{centre} and on the expensive side. \\ \cdashline{1-2}

Retrieved 1 &  \textcolor{myOrange}{kokai} means \textcolor{myBlue}{open to} the public or laid \textcolor{myBlue}{open}      \\
\quad \small{+ ASR} & cook eye means open to the public all laid open  \\

Retrieved 2&   the town of beauharnois was the major \textcolor{myGreen}{centre}    \\
\quad \small{+ ASR} & the town of bo hanwa was the major center    \\
Retrieved 3               &  the gate is \textcolor{myBlue}{open} at eleven    \\
\quad \small{+ ASR} & the gate is open at 11   \\
\bottomrule
\end{tabular}
}
\caption{In-context examples retrieved based on phoneme-level similarity.
For each target utterance (top row), we retrieve three ($g$, $e$) example pairs from the database using phonetic similarity between the target and $g$.  Colored segments highlight phonetically similar phrases between the target and retrieved examples.}
\label{tab:retreive_example}
\end{table*}



\subsection{Retrieved In-Context Example}
\label{app:EPA_retreive}
In Section~\ref{sec:method_step1}, we retrieved in-context examples based on phoneme-level similarity. Table~\ref{tab:retreive_example}, we present several representative examples to illustrate this retrieval process, showing how phonetically similar phrases (highlighted in color) are matched between the target and retrieved utterances. This demonstrates that the retrieval mechanism effectively captures pronunciation-level patterns relevant to ASR-style errors.


\subsection{Additional Examples of ASR-style Errors}

\begin{table*}[h!]
\small
\resizebox{\textwidth}{!}{
\begin{tabular}{p{0.5\textwidth}p{0.5\textwidth}}


\toprule
\textbf{Original} & \textbf{Augmented} \\\hline

- no, i just need to make sure it's \textcolor{myOrange}{cheap}. oh, and i need parking  
& - no, i just need to make sure it's \textcolor{myOrange}{sheep}. oh, and i need parking. \\

- i am departing from \textcolor{myBlue}{marion}  - 
& - i am departing from \textcolor{myBlue}{maryland} \\

- no, i don't need  \textcolor{myOrange}{anything} else right now. thank you for your  \textcolor{myOrange}{assistance}. good  \textcolor{myOrange}{bye}.  
& - no, i don't need \textcolor{myOrange}{anyone} else right now. thank you for your \textcolor{myOrange}{persistence}. good \textcolor{myOrange}{buy}. \\

- i would like to go to \textcolor{myBlue}{sandy} please.  
& - i would like to go to \textcolor{myBlue}{cindy} please. \\

- i would like to keep it in the \textcolor{myBlue}{moderate} range, please.  
& - i would like to keep it in the \textcolor{myBlue}{mod rain}, please \\

- could i get the \textcolor{myOrange}{address} for it? i would also like an \textcolor{myBlue}{expensive} place to eat around it.  
& - could i get the \textcolor{myOrange}{actress} for it? i would also like an \textcolor{myBlue}{extensive} place to eat around it. \\

- i need to take a \textcolor{myBlue}{train} out of garrett, i will be leaving town on wednesday
& - i need to make a \textcolor{myBlue}{plane} out of garrett, i will be \textcolor{myOrange}{weaving} town on wednesday. \\

- do you have any \textcolor{myBlue}{indian} restaurants in the \textcolor{myBlue}{south} in a different \textcolor{myOrange}{price} range?  
& - do you have any \textcolor{myBlue}{indonesian} restaurants in the \textcolor{myBlue}{south} in a different \textcolor{myOrange}{prize} range \\

- nope, same \textcolor{myOrange}{people}.  
&-  nope, same \textcolor{myOrange}{pupil}. \\

- i'm looking for a \textcolor{myBlue}{college} type \textcolor{myBlue}{attraction}  
& - i'm looking for a \textcolor{myBlue}{knowledge}-type \textcolor{myBlue}{action}. \\

- yes, please \textcolor{myOrange}{book} me a \textcolor{myOrange}{room} for friday
& - yes, please \textcolor{myOrange}{cook} me a \textcolor{myOrange}{broom} for friday \\

- yes, could you please email or \textcolor{myOrange}{fax} me the \textcolor{myOrange}{fare} amount, as well as the reference number?  
& - yes, could you please email or \textcolor{myOrange}{text} me the \textcolor{myOrange}{fair} amount, as well as thereference code? \\

- \textcolor{myBlue}{ois el shaddai} a \textcolor{myOrange}{guest house} or \textcolor{myBlue}{hotel}?  
& \textcolor{myBlue}{iz let shadai} a \textcolor{myOrange}{gest house} or \textcolor{myBlue}{motel}? \\

- great! i also need a \textcolor{myBlue}{train} from \textcolor{myBlue}{mount pleasant} to \textcolor{myBlue}{sabattus}, please.  
& - great! i also need a \textcolor{myBlue}{strain} from \textcolor{myBlue}{mount pleasant} to \textcolor{myBlue}{suspicious}, please \\

- yes, can you help me find a \textcolor{myBlue}{train} that can take me from \textcolor{myBlue}{lovelock} to \textcolor{myBlue}{abbot}?  
& - yes, can you help me find a \textcolor{myBlue}{plane} that can take me from \textcolor{myBlue}{love lock} to \textcolor{myBlue}{rabbit}? \\


\bottomrule
\end{tabular}
}
\caption{
Examples of augmented utterances generated by injecting phoneme-level ASR-style errors. For each original utterance (left), the corresponding augmented version (right) includes substitutions that mimic realistic ASR recognition mistakes.  
\textcolor{myBlue}{Blue-colored phrases} indicate changes in keywords that are used as slot values in DST, while \textcolor{myOrange}{orange-colored phrases} represent overall ASR-style errors.
}

\label{tab:table_asr_examples}
\end{table*}



\label{app:asr-examples}

Table~\ref{tab:table_asr_examples} presents additional examples of ASR-style errors generated by our EPA method, including both general and keyword-specific transformations.



\section{Experimental Setup}
\label{app:exp_setting}
\subsection{Details of the ASR Environment}
\begin{itemize}
    \item Low-acc ASR environment: Whisper-base model (74M)\cite{whisper} is used for transcription. WER on LibriSpeech.test-clean is 0.05.
    \item Noisy audio environment: Incorporated authentic cafe and traffic noise from \url{https://freesound.org/} with a 10 to 20 Signal-to-Noise Ratio (SNR) and transcribed it using the Whisper large model.
    \item Paraphrased environment: When recording the audio, the text was paraphrased to resemble more natural, real-life spoken language\cite{dstc11}.
    \item High-acc ASR environment: Whisper-large model (1550M) is used for transcription. WER on LibriSpeech.test-clean is 0.027.
\end{itemize}




\subsection{Comparison Methods} 
\label{sec:appendix-experiment-compared-methods}
\begin{itemize}
    \item AEDA \cite{aeda}: We randomly inserted punctuation marks, effectively maintaining the original word order.
    \item EDA \cite{eda}: We augmented data by applying edit-based technique that implements four rule-based modifications—synonym replacement, random insertion, swapping, and deletion.
    \item Back Translation \cite{BT}: We translated original texts to error texts and then back to the original texts for generating syntactic variations during the process. We use English to German \footnote{\url{facebook/wmt19-en-de}} and German to English\footnote{\url{facebook/wmt19-de-en}} models as translator.
    
    
    \item TTS-ASR : We used Tacotron2 \cite{taco} for the TTS model to synthesize the audio and use Whisper-base \cite{whisper} as an ASR model to simulate the ASR errors.

    
    \item ASR translation: We employed a sequence-to-sequence structure to translate clean text into ASR-errored text. Our training set comprised 300 hours of paired clean and ASR-errored text. We fine-tuned the model based on the T5-base architecture \cite{t5}, using the loss function defined in equation~\ref{eq:label_loss_translation}. The loss function is as follows:
    \begin{equation}\label{eq:label_loss_translation}
      L = -\sum_{i=1}^{I} \log P(e_i|g_i).
    \end{equation}

\end{itemize}



\subsection{Training Details}
\label{sec:app-training}
 In training models, we used T5-base \cite{t5} as the backbone model and instructed the model to generate the $B_t$ by given $D_t$ in sequence to sequence manner, as in \cite{pptod} and the loss function is 
    \begin{equation}\label{eq:label_loss}
      L = -\sum_{t=1}^{T} \log P(B_t|\text{Inst}, D_t).
    \end{equation}
We set the learning rate as 4e-5 and used the AdamW \cite{loshchilov2017decoupled} optimizer. One GeForce RTX 3090 is used for training and the batch size is 16. Trained until reaching the max patient, which is 3.




\section{Further Experiments}
\label{app:exp_main}



\subsection{Statistical Significance Analysis}
\begin{table}[t]
\centering
\resizebox{\columnwidth}{!}{%
\begin{tabular}{l|cc|cc|cc|cc}
\hline
\multicolumn{1}{c|}{\multirow{2}{*}{Method}} &
  \multicolumn{2}{c|}{Low-acc ASR} &
  \multicolumn{2}{c|}{Noised Aud.} &
  \multicolumn{2}{c|}{Paraphrased} &
  \multicolumn{2}{c}{High-acc ASR} \\ \cline{2-9} 
\multicolumn{1}{c|}{} & JGA   & N-acc & JGA   & N-acc & JGA   & N-acc & JGA   & N-acc \\ \hline
Baseline              & --    & --    & --    & --    & --    & --    & --    & --    \\ \hline
AEDA                  & ns    & ns    & ns    & ns    & ns    & ns    & ns    & ns    \\
EDA                   & *     & *     & **    & **    & ns    & **    & **    & **    \\
BT                    & **    & **    & **    & **    & *     & **    & **    & *     \\
TTS-ASR               & ns    & ns    & ns    & ns    & ns    & ns    & ns    & ns    \\
ASR trans.                   & ns    & *     & ns    & ns    & ns    & *     & *     & *     \\
EPA (GPT-3.5)         & ***   & **    & ***   & ***   & **    & **    & *     & **    \\
\hline
\end{tabular}%
}
\caption{Statistical significance results compared to the Baseline using paired \textit{t}-tests across three random seeds. Stars indicate significance levels: * for $p<0.05$, ** for $p<0.01$, *** for $p<0.001$, and \textit{ns} for non-significant differences.}
\label{tab:significance}
\end{table}

To assess the reliability of our results, we conducted paired \textit{t}-tests between each method and the Baseline to determine whether the observed performance improvements are statistically significant. We report 95\% confidence intervals to reflect performance variability. Statistical significance is denoted using asterisks: $\ast$ for $p < 0.05$, $\ast\ast$ for $p < 0.01$, and $\ast\ast\ast$ for $p < 0.001$.

As shown in Table~\ref{tab:significance}, EPA (GPT-3.5) achieves statistically significant gains in nearly all evaluation settings, particularly under low-accuracy and noised ASR conditions. These results confirm that the improvements brought by our method are both consistent and statistically reliable.



\subsection{Baseline Performance Comparison with Clean Text }
For comparison, we report the baseline performance on an error-free, clean test dataset. Please note that DSTC11 \cite{dstc11} does not provide a text script for the test dataset, so we are manually cleaning 50 dialogues to ensure they are error-free. In the experiment, the baseline model achieved a JGA score of 45.2 \% and an N-ACC score of 86.5 \% in an ASR error-free environment. Compared to the JGA, which is 34.8 \%, and N-ACC, which is 52.07 \%, in the ASR-errored environment (High-acc ASR model environment), this discrepancy highlights the significant impact of ASR errors on performance degradation.



 
\subsection{Effect of Keyword Augmentation}
\begin{table}[t]
\centering
\resizebox{\columnwidth}{!}{%
\begin{tabular}{l|cc|cc|cc}
\hline
\multirow{2}{*}{Method} &
\multicolumn{2}{c|}{Weak-ASR } &
\multicolumn{2}{c|}{Noised Audio} &
\multicolumn{2}{c}{Strong  ASR} \\ \cline{2-7}
 & JGA & N-acc & JGA & N-acc & JGA & N-acc \\ \hline
Baseline & 29.88 & 45.76 & 29.70 & 46.77 & 34.87 & 52.07 \\
+ Overall Error (Step 1) & 31.31 & 50.67 & 31.13 & 52.29 & 35.40 & 55.80 \\
+ Keyword Error (Step 2) & 31.66 & 49.88 & 31.61 & 51.02 & 36.12 & 54.23 \\
EPA (Step1+2) & \textbf{32.39} & \textbf{51.12} & \textbf{32.24} & \textbf{52.70} & \textbf{36.61} & \textbf{55.87} \\ \hline
\end{tabular}%
}
\caption{Ablation on overall and keyword-specific augmentation steps. Combining both (EPA) yields the highest robustness across ASR conditions.}
\label{tab:keyword_ablation}
\end{table}

To further validate the contribution of each component in our method, we conducted an ablation study separating Step 1 (overall error augmentation) and Step 2 (keyword-specific error augmentation), using GPT-3.5 as the generator. Both augmentation strategies individually improved performance over the baseline. 
Step~1 yielded consistent gains across all ASR settings, while Step~2 notably enhanced slot-value accuracy. When combined (EPA, ours), the model achieved the highest robustness in all evaluation conditions, confirming that the two strategies are complementary.




\subsection{Experiments with a Different Baseline}
\begin{table}[t]
\centering
\resizebox{\columnwidth}{!}{%
\begin{tabular}{l|cc|cc|cc}
\hline
\multicolumn{1}{c|}{\multirow{2}{*}{Method}} &
  \multicolumn{2}{c|}{Low-acc ASR} &
  \multicolumn{2}{c|}{Noised Aud.} &
  \multicolumn{2}{c|}{Paraphrased} \\\cline{2-7} 
\multicolumn{1}{c|}{} & JGA   & N-acc & JGA   & N-acc & JGA   & N-acc  \\ \hline
Baseline              &29.9 & 45.82 & 27.25 & 41.81 & 25.81 & 44.51 \\ 
% \quad + EPA (Step1)               & 12.68 & 81.85 & 12.94 & 81.97 & 13.05 & 77.13 \\ 
\quad + EPA             &30.63 & 48.54 & 27.5 & 46.52 & 27.65 & 47.96 \\ \hline
\end{tabular}%
}
\caption{Experiment with GPT-2 model as baseline.}
\label{tab:gpt2}
\end{table}





% \begin{table}[t]
% \centering
% \resizebox{\columnwidth}{!}{%
% \begin{tabular}{lcccc}
% \hline
% & \textbf{Low-acc ASR} & \textbf{Noised Aud.} & \textbf{Paraphrased} & \textbf{High-acc ASR} \\
% \hline
% Baseline & 12.68 & 12.45 & 13.05 & 12.82 \\
% + EPA    & 15.90 & 15.29 & 15.07 & 16.62 \\
% \hline
% \end{tabular}
% }
% \caption{Experiment with GPT-2 model as baseline. JGA is reported.}
% \label{tab:gpt2}
% \end{table}



In the main experiments, we use T5-base as the backbone model. To assess the generalizability of our approach, we additionally conduct experiments using a GPT-2~\cite{gpt2} model, as shown in Table~\ref{tab:gpt2}. The results show a consistent trend with those of T5-base, demonstrating that our method is effective across different backbone architectures.



\subsection{Generalizability to Other Tasks}
\label{app:generalizability}
\begin{table}[t]
\centering
\resizebox{\columnwidth}{!}{%
\begin{tabular}{lccc}
\hline
\textbf{Method} & \textbf{Low-acc ASR} & \textbf{Noised Audio} & \textbf{High-acc ASR} \\
\hline
Baseline             & 56.29 & 60.50 & 60.02 \\
+ OPT 6.7B           & 59.64 & 62.84 & 62.47 \\
+ LLaMA 7B           & 58.29 & 60.53 & 60.46 \\
+ GPT-3.5 (125B)     & 59.39 & 62.62 & 62.16 \\
\hline
\end{tabular}
}
\caption{NER results on the ASAPP/SLUE dataset under different ASR conditions. We reported the F1 score.}
\label{tab:ner_results}
\end{table}
\begin{table}[t]
\centering
\resizebox{\columnwidth}{!}{%
\begin{tabular}{lccc}
\hline
\textbf{Method} & \textbf{Low-acc ASR} & \textbf{Noised Audio} & \textbf{High-acc ASR} \\
\hline
Baseline             & 55.25 & 63.23 & 64.32 \\
+ OPT 6.7B           & 56.97 & 64.94 & 66.10 \\
+ LLaMA 7B           & 57.24 & 64.26 & 65.56 \\
+ GPT-3.5 (125B)     & 58.78 & 66.16 & 67.49 \\
\hline
\end{tabular}
}
\caption{SLU results on the SLURP dataset under different ASR conditions. We reported F1 score.}
\label{tab:slu_results}
\end{table}

To evaluate the generalizability of our approach beyond the DST domain, we extended our experiments to two additional spoken language understanding tasks: Named Entity Recognition (NER) and Spoken Language Understanding (SLU). We applied our EPA methodology under three ASR conditions—low-accuracy ASR, noised audio, and high-accuracy ASR—using the same experimental setup as in the DSTC11 evaluation. We used asapp/slue dataset for NER task\cite{asapp}, and  SLURP dataset \cite{slurp} for SLU task.

The results, shown in Table~\ref{tab:ner_results} and Table~\ref{tab:slu_results}, demonstrate that our method consistently improves performance across all ASR conditions for both NER and SLU tasks. Notably, the gains are especially prominent under low-accuracy and noisy conditions, confirming that our approach is broadly applicable to other tasks.



\subsection{Additional Fine-grained Metrics}
\label{appendix:finegrained}
\begin{table}[h]
\centering
\resizebox{\columnwidth}{!}{%
\begin{tabular}{lcccc}
\hline
\textbf{Method} & \textbf{Precision} & \textbf{Recall} & \textbf{F1} & \textbf{Slot Accuracy} \\
\hline
Baseline         & 50.5 & 50.1 & 50.3 & 92.8 \\
+ TTS-ASR          & 49.5 & 49.6 & 49.5 & 93.4 \\
+ ASR-Translation  & 51.8 & 51.7 & 51.7 & 93.1 \\
+ EPA (GPT3.5)     & 55.2 & 54.8 & 55.0 & 93.7 \\
\hline
\end{tabular}
}
\caption{Fine-grained DST metrics under Low-accuracy ASR setting.}
\label{tab:fine_1}
\end{table}


\begin{table}[h]
\centering
\resizebox{\columnwidth}{!}{%
\begin{tabular}{lcccc}
\hline
\textbf{Method} & \textbf{Precision} & \textbf{Recall} & \textbf{F1} & \textbf{Slot Accuracy} \\
\hline
Baseline         & 52.0 & 51.3 & 51.6 & 92.5 \\
+ TTS-ASR          & 53.7 & 53.1 & 53.4 & 92.8 \\
+ ASR-Translation  & 53.8 & 53.0 & 53.4 & 92.5 \\
+ EPA (GPT3.5)     & 56.1 & 55.3 & 55.7 & 93.2 \\
\hline
\end{tabular}
}
\caption{Fine-grained DST metrics under Noised ASR setting.}
\label{tab:fine_2}
\end{table}

To supplement the main results focusing on JGA and N-Acc, we report additional fine-grained metrics—Precision, Recall, F1, and Slot Accuracy—under two ASR corruption settings: Low-accuracy ASR and Noised ASR. These metrics provide a more comprehensive view of model behavior in diverse error conditions in table~\ref{tab:fine_1} and ~\ref{tab:fine_2}.







\subsection{Results with Different Random Seeds}
% seed2
\begin{table}[t]
\centering
\resizebox{\columnwidth}{!}{%
\begin{tabular}{l|cc|cc|cc|cc}
\hline
\multicolumn{1}{c|}{\multirow{2}{*}{Method}} &
  \multicolumn{2}{c|}{Low-acc ASR} &
  \multicolumn{2}{c|}{Noised Aud.} &
  \multicolumn{2}{c|}{Paraphrased} &
  \multicolumn{2}{c}{High-acc ASR} \\ \cline{2-9} 
\multicolumn{1}{c|}{} & JGA   & N-acc & JGA   & N-acc & JGA   & N-acc & JGA   & N-acc \\ \hline
Baseline              & 30.05 & 46.48 & 29.80 & 47.49 & 29.08 & 48.85 & 34.73 & 52.31 \\ \hline
AEDA                  & 29.99 & 46.53 & 29.80 & 47.60 & 28.95 & 49.02 & 34.94 & 52.40 \\
EDA                   & 29.16 & 47.67 & 28.93 & 49.50 & 28.12 & 50.10 & 33.74 & 53.66 \\
BT                    & 31.54 & 49.21 & 31.43 & 51.25 & 29.90 & 51.60 & 36.25 & 54.85 \\
TTS-ASR               & 30.09 & 46.16 & 30.32 & 47.74 & 29.26 & 49.31 & 35.28 & 51.95 \\
ASR trans.                        & 29.98 & 47.70 & 29.95 & 48.49 & 29.72 & 50.43 & 34.82 & 53.35 \\
% EPA (Llama 7B)        & 31.54 & 51.12 & 31.55 & 52.27 & 30.10 & 53.55 & 36.22 & 55.49 \\
% EPA (Opt 6.7B)        & 31.82 & 50.73 & 32.03 & 51.92 & 29.57 & 52.49 & 37.05 & 55.78 \\
EPA (GPT3.5)          & 32.56 & 51.62 & 32.27 & 53.48 & 31.10 & 53.91 & 36.65 & 56.14 \\ \hline \hline
Baseline              & 29.82 & 45.63 & 29.75 & 46.20 & 28.55 & 48.75 & 34.83 & 51.44 \\ \hline
AEDA                  & 29.77 & 46.29 & 29.69 & 47.40 & 29.26 & 48.93 & 34.93 & 52.26 \\
EDA                   & 29.27 & 47.48 & 28.57 & 49.41 & 28.20 & 49.99 & 33.59 & 53.68 \\
BT                    & 31.66 & 49.16 & 31.05 & 50.81 & 29.88 & 51.75 & 36.19 & 54.73 \\
TTS-ASR               & 29.75 & 45.99 & 29.80 & 46.87 & 29.10 & 48.12 & 34.97 & 51.49 \\
ASR trans.                        & 30.86 & 47.47 & 30.52 & 48.53 & 29.63 & 50.11 & 35.73 & 53.59 \\
% EPA (Llama 7B)        & 31.76 & 50.63 & 31.95 & 51.93 & 30.83 & 52.82 & 36.52 & 55.29 \\
% EPA (Opt 6.7B)        & 31.12 & 49.46 & 30.90 & 50.85 & 29.92 & 51.65 & 35.81 & 54.78 \\
EPA (GPT3.5)          & 32.33 & 50.74 & 32.41 & 52.28 & 31.02 & 53.15 & 36.75 & 55.82 \\ \hline\hline
Baseline              & 29.77 & 45.18 & 29.56 & 46.61 & 29.12 & 48.76 & 35.05 & 52.47 \\ \hline
AEDA                  & 29.94 & 46.55 & 29.72 & 47.45 & 29.16 & 48.62 & 34.94 & 52.29 \\
EDA                   & 29.22 & 47.79 & 28.59 & 49.61 & 27.92 & 49.88 & 33.71 & 53.99 \\
BT                    & 31.86 & 49.13 & 31.31 & 50.87 & 29.91 & 51.85 & 36.38 & 54.85 \\
TTS-ASR               & 29.98 & 46.86 & 29.84 & 47.51 & 28.88 & 49.20 & 34.97 & 52.66 \\
ASR trans.                        & 30.37 & 47.78 & 29.94 & 48.32 & 29.26 & 50.60 & 35.20 & 54.05 \\
% EPA (Llama 7B)        & 32.09 & 50.36 & 31.93 & 51.68 & 30.63 & 53.00 & 36.50 & 55.31 \\
% EPA (Opt 6.7B)        & 31.09 & 49.73 & 31.06 & 51.02 & 29.95 & 51.79 & 35.88 & 54.69 \\
EPA (GPT3.5)          & 32.27 & 51.01 & 32.04 & 52.33 & 30.72 & 52.96 & 36.42 & 55.66 \\ \hline
\end{tabular}%
}
\caption{Experiment result with different seeds.}
\label{tab:seeds}
\end{table}


Table~\ref{tab:seeds} reports the results of our main experiments (Table~\ref{tab:main}) repeated with three different random seeds, demonstrating the consistency of the observed trends.







\section{Details of Quality Evaluation}
\label{app:quality}
\subsection{About Metric}

As described in Section~\ref{sec:quality}, we use a phonetic similarity metric to evaluate pronunciation-level consistency between the original and augmented text. Specifically, we compute the normalized phoneme edit distance, which quantifies the minimal number of phoneme-level operations required to transform one utterance into another. A higher score indicates greater phonetic similarity.
We used the \texttt{eng-to-ipa} library\footnote{\url{https://pypi.org/project/eng-to-ipa/}} for phoneme conversion in our implementation, as shown in the code snippet below.

\begin{lstlisting}[style=pythonstyle]
def phonetic_similarity(original_text, augmented_text):
    
    original_ipa  = to_phoneme(original_text)
    augmented_ipa = to_phoneme(augmented_text)
    
    edit_distance = nltk.edit_distance(original_ipa, augmented_ipa)
    
    # Normalize the edit distance
    max_length = max(len(original_ipa), len(augmented_ipa))
    normalized_distance = float(edit_distance) / float(max_length)
    
    # Convert to similarity score
    similarity_score = 1 - normalized_distance
    return similarity_score
\end{lstlisting}


\subsection{Human Evaluation Details}

To assess the plausibility of the generated ASR-style errors, we conducted a human evaluation involving two graduate students. Each participant was asked to rate whether a given sentence transformation could plausibly be attributed to an ASR error, using a 4-point Likert scale:

\begin{itemize}
    \item \textbf{1 – Not at all}: The change is unlikely to be due to an ASR error. It appears to stem from other factors such as meaning alteration or stylistic variation.
    \item \textbf{2 – Unlikely}: The transformation is probably not caused by an ASR error.
    \item \textbf{3 – Somewhat likely}: The transformation may plausibly be caused by an ASR error.
    \item \textbf{4 – Very likely}: The transformation clearly appears to result from an ASR error.
\end{itemize}

Each sentence pair (original and transformed) was rated independently by both annotators. Inter-rater agreement and average scores are reported in Section~\ref{sec:quality}. The distribution of Likert scores for each annotator is as follows : Annotator 1 assigned 5\% of scores as 1, 7\% as 2, 17\% as 3, and 71\% as 4. Annotator 2 assigned 3\% of scores as 1, 24\% as 2, 53\% as 3, and 20\% as 4.








% \textbf{Step 1 Prompt Example} \\
% \small{\texttt{Generate ASR error augmented text with similar pronunciation but different words based on the given gold text examples.\\
% Apply character and word substitutions, additions, or deletions while maintaining the overall pronunciation and context. \\
% Error rate should be high \\
% \{Example 1\}\\
% Original: they have a single naupliar eye  \\
% ASR-errored: they have a single nor pure eye \\
% \{Example 2\}\\
% Original: i must have saint louis then huzza  \\
% ASR-errored: i must have st louis then hazard \\
% \{Example 3\}\\
% Original: i wonder uncle did not have her come  \\
% ASR-errored: i wonder uncle did not have a problem \\
% \{Example 4\}\\
% Original: i don't want to have to pay for parking. \\
% ASR-errored: \\ \\ \\  
% }}
% \normalsize\textbf{Example2}
% \\
% \noindent\small{\texttt{{Generate ASR error augmented text with similar pronunciation but different words based on the given gold text examples.\\
% Apply character and word substitutions, additions, or deletions while maintaining the overall pronunciation and context. \\
% Error rate should be high \\
% \{Example 1\}\\
% Original: anatomical information for the southern fur seals arctocephalus species is scant  \\
% ASR-errored: anatomical information for the southern thursills aquatopheas species is kept \\
% \{Example 2\}\\
% Original: and how do you find yourself my poor lad  \\
% ASR-errored: and thank you for yourself my poor lad \\
% \{Example 3\}\\
% Original: the town of beauharnois was the major centre  \\
% ASR-errored: the town of bo hanwa was the major center \\
% \{Example 4\}\\
% Original: i'm looking for a guesthouse in the centre of town. can you find one for me? \\
% ASR-errored: \\
% }}}
% \\
% \noindent
% \normalsize\textbf{Example3}
% \\
% \noindent\small{\texttt{{Generate ASR error augmented text with similar pronunciation but different words based on the given gold text examples.\\
% Apply character and word substitutions, additions, or deletions while maintaining the overall pronunciation and context. \\
% Error rate should be high \\
% \{Example 1\}\\
% Original: i need fifty ten foot long segments of wire  \\
% ASR-errored: i need fifty ten foot long signals of my life \\
% \{Example 2\}\\
% Original: i d go into a trance and see if i could locate my chum  \\
% ASR-errored: i d go into a trance and see if i could look at what s on \\
% \{Example 3\}\\
% Original: i go there every year and my daughter visits me  \\
% ASR-errored: i m going to go there every year and my daughter visits me \\
% \{Example 4\}\\
% Original: i need to go from ashland to peterbourgh saturday. \\
% ASR-errored: \\
% }}}



% \section{Experiment with Different Random Seed}
% \label{sec:random_seed}
% \normalsize{
% In  Table~\ref{tab:seeds}, we report the experiment result with different random seed of our main table(Table~\ref{tab:main}).}

% \input{table_seed}


% Uniq words, N. errored 제외한 column, 표3개 그리기
% SOL


% In this section, we describe the method employed for calculating Shannon Entropy to measure the diversity of words in a given set of sentences. The Shannon Entropy ($H$) is calculated based on the probability distribution of words within the text.

% Let $N$ be the total number of words in the corpus, and $p_i$ be the probability of occurrence of the $i$-th unique word in the corpus. The Shannon Entropy is then defined as:

% \[
% H = -\sum_{i=1}^{N} p_i \cdot \log_2(p_i)
% \]

% To apply this method to our dataset, we follow the steps outlined below:

% \begin{enumerate}
%     \item \textbf{Word Frequency Calculation:}
%     \begin{itemize}
%         \item Count the occurrences of each unique word in the entire dataset.
%     \end{itemize}
    
%     \item \textbf{Probability Calculation:}
%     \begin{itemize}
%         \item Calculate the probability ($p_i$) of each unique word by dividing its frequency by the total number of words in the dataset.
%     \end{itemize}
    
%     \item \textbf{Shannon Entropy Calculation:}
%     \begin{itemize}
%         \item Apply the Shannon Entropy formula using the calculated probabilities.
%     \end{itemize}
% \end{enumerate}

% \textbf{Results:}
% The obtained Shannon Entropy value provides a measure of the diversity of words in the given set of sentences.

% The detailed calculations, including the word frequency distribution and probability assignments, are presented in Appendix A for reference.



% \section{Prompt Example}
% \label{sec:prompt_example}


\normalsize


